\begin{figure}
    \center
    \begin{tikzpicture}[ % FISSA IL RIFERIMENTO IN MODO CORRETTO, ALTRIMENTRI --> (X,Z,Y)
            x={(0.866cm,-0.5cm)},   % cos30, -sin30
            y={(0.866cm,0.5cm)},    % cos30,  sin30
            z={(0cm,0.8cm)},
            scale = 1.,
            name path global/.style={},
            ]
        % Parametri
        \def \xM {2}
        \def \yM {2}
        \def \R {1.5}
        \def \h {4}
        \coordinate (M) at ({\xM},{\yM},0); % centro della base inferiore
        \coordinate (N) at ({\xM},{\yM},\h); % centro della base superiore
        
        % Assi cartesiani
        \draw[->] (0,0,0) -- ({\xM + 1},0,0) node[below left] {$x$};
        \draw[->] (0,0,0) -- (0,{\yM + 1},0) node[above left] {$y$}; % below/above
        \draw[->] (0,0,0) -- (0,0,{\h + 1}) node[left] {$z$};

        % basi
        % base inferiore come path nominato
        \path[name path=BaseCircle] (M) circle [radius={\R}];
        \draw[name path=BaseCircleDraw] (M) circle [radius={\R}];
        \draw (N) circle [radius={\R}];
        % altezze

        % punti estremi sinistra/destra della base inferiore
        %\coordinate (LeftBase)  at ({\xM - \R},{\yM},0);
        %\coordinate (RightBase) at ({\xM + \R},{\yM},0);

        % generatrici verticali dai punti estremi
        %\draw (LeftBase)  -- ++(0,0,\h); % il ++ aggiunge il vettore al vettore precedente (coordinate relative)
        %\draw (RightBase) -- ++(0,0,\h);

        % retta nel piano xy (esempio)
        \path[name path=MyLine] (0,0,0) -- (4,4,0);
        %\draw[thick] (0,1,0) -- (5,3,0);

        % calcolo intersezioni

        \path[name intersections={of=BaseCircle and MyLine, by={I1,I2}}];
        % Salva I1 e I2 come coordinate vere
        \coordinate (C1) at (I1);
        \coordinate (C2) at (I2);

        \draw (I1) -- ++(0,0,\h);
        \draw (I2) -- ++(0,0,\h);
        % punti di intersezione evidenziati
        %\fill (I1) circle (2pt);
        %\fill (I2) circle (2pt);

        \draw[->] (N) -- ++(0,0,1) node[right] {$\mathbf{ M }$};


        % Calcola centro e raggio della semicirconferenza
        \coordinate (CenterArc) at ($(N)+(0,0,1)$);
        % Calcola i punti di partenza e arrivo
        \coordinate (StartArc) at ($(C2)+(0,0,1)$);
        \coordinate (EndArc) at ($(C1)+(0,0,1)$);

        % Calcola angoli rispetto al centro
        \pgfmathanglebetweenpoints{\pgfpointanchor{CenterArc}{center}}{\pgfpointanchor{StartArc}{center}}
        \let\StartAngle\pgfmathresult
        \pgfmathanglebetweenpoints{\pgfpointanchor{CenterArc}{center}}{\pgfpointanchor{EndArc}{center}}
        \let\EndAngle\pgfmathresult
        % Calcola raggio
        \pgfmathsetmacro{\ArcRadius}{veclen((\xM+\R)-(\xM), (\yM)-(\yM))}

        % Disegna la semicirconferenza
        \draw[thick,->] (StartArc) arc[start angle=\StartAngle, end angle=405, radius=\ArcRadius] node[above right] {$ \mathbf{ K }_m $};


    \end{tikzpicture}
    \caption{Cilindro con vettore magnetizzazione}
    \label{figure:_cilindro_con_magnetizzazione}
\end{figure}